Para la ejecución de los experimentos se utilizó una laptop ASUS Zenbook 14. Entre sus componentes de hardware cuenta con un, Intel Core Ultra 7 155H con 15 GB de memoria RAM. En cuanto a software, el sistema operativo empleado fue Fedora Linux 44 (Workstation Edition) de 64 bits. El código fuente fue compilado utilizando GCC (g++) en su versión 16.2.1, aplicando la "flag" de optimización \texttt{-O3}.

Respecto a los \textit{datasets}. Para los algoritmos de ordenamiento, se crearon arreglos unidimensionales de números enteros con tamaños $N \in \{10^1, 10^3, 10^5, 10^7\}$, evaluando distribuciones ascendentes, descendentes y aleatorias, bajo los dominios D1 (rango: $[0,9]$) y D7 (rango $[0,10^{7}$). En cuanto a la multiplicación de matrices cuadradas, se generaron matrices de dimensión $N \in \{2^4, 2^6, 2^8, 2^{10}\}$, abarcando tipos dispersos, diagonales y densos, con coeficientes en los dominios D0 (0 o 1) y D10 (rango: $[0,9]$). Para cada combinación de parámetros se midieron tres muestras (a, b y c). 
